\lamina{img/il_fractal.jpg}
\capitulo{CAPÍTULO 24}{LA REALIDAD FRACTAL}{Sobre la imaginación como creación, continuación de \emph{El diapasón invisible}}{}
\noindent \capital{C}{asi} nadie se hace esta pregunta, pero aparece sola en cuanto alguien lleva años imaginando un dragón: no qué veo yo cuando pienso en él, sino qué ve él.

Desde dentro de su cueva, el dragón no tiene manera de asomarse al cuarto donde lo estoy pensando. Puede mirar el techo de roca todo lo que quiera: no me verá. Y sin embargo yo estoy ahí, decidiendo si esta noche vuela.

Es la única relación de este tipo que puedo ver entera, porque en ella me toca la parte de arriba. Y en cuanto la veo entera, se vuelve incómoda por el otro lado: si yo soy eso para el dragón, ¿cómo sé que no hay, sobre mi propio techo, alguien en la misma posición respecto a mí?

\seccion{Un cuento chino muy viejo}
\noindent El taoísmo gira en torno a una sola idea: el Tao, el fondo del que sale todo sin que él salga de nada. En el capítulo 5 contamos la historia que el Zhuangzi dedica a Hun Dun, el emperador del centro: una plenitud sin rasgos, sin los siete agujeros con los que los demás ven, oyen, respiran y comen.

Recordemos: dos reyes vecinos, agradecidos por su hospitalidad, deciden pagársela abriéndole un agujero cada día. Al séptimo, Hun Dun ha dejado de existir.

Es fácil leer esa muerte como el nacimiento de la conciencia normal: la plenitud sin rasgos se rompe en percepciones, y donde había un fondo continuo aparece un mundo con bordes. Ese fondo sin rasgos ya tenía nombre en este libro: el reservorio, el océano del que salen las olas y al que vuelven.

Pero hay algo que esa lectura no explica. Si Hun Dun fuera el Tao mismo, no podría morir: lo anterior a todo no tiene por dónde acabarse. Que muera dice, más bien, que no es la fuente. Es la cara que la fuente enseña hacia fuera. Morir no destruye el origen. Solo destruye, para quien acaba de abrir los ojos, la superficie sin rasgos —y a cambio le entrega un mundo con bordes.

\begin{fisica}
\lbl{En física esto se llama:} la interfaz no es la fuente.

\lbl{En la vida diaria es como:} la puerta de una casa. Puedes tirarla abajo sin que la casa pierda ni una habitación. Lo que pierdes es la forma de entrar.
\end{fisica}
\seccion{Tres verbos y nada más}
\noindent Si esto vale para Hun Dun, vale también para el reservorio: no es la última palabra, sino un puñado de operaciones que algo, más atrás, nos ofrece. Y son pocas. Aparece un borde donde no lo había. Lo que está dentro del borde cambia de forma sin dejar de tener borde. El borde se disuelve y la información vuelve al fondo. Nacer, vivir y morir no son tres misterios distintos: son los nombres caseros de esos tres verbos, los mismos que el capítulo 8 tradujo a instanciar, ejecutar y liberar.

\begin{fisica}
\lbl{En física esto se llama:} conservación de la información.

\lbl{En la vida diaria es como:} el ciclo del agua. La nube, la lluvia y el río no son tres cosas: son tres formas de la misma agua.
\end{fisica}
Pero si todo sale de la misma agua, ¿por qué nadie se repite? Dos olas son la misma agua y no son la misma ola: lo que no se repite es la forma en que se enrollaron, el momento exacto en que cuajaron. Cada pliegue, por pequeño que sea, es un dentro completo, visto desde su propio centro. No es un trozo de otro.

\begin{fisica}
\lbl{En física esto se llama:} cada uno cristaliza a su manera.

\lbl{En la vida diaria es como:} dos olas que son la misma agua y no son la misma ola.
\end{fisica}
\seccion{¿Por qué soy más real que mi dragón?}
\noindent La única interioridad a la que tengo acceso directo es la mía. La del vecino, la del perro, se infieren —con buenas razones, pero se infieren. Y un dragón que alguien lleva cincuenta años imaginando también se comporta, también sorprende, también se niega a veces a volar cuando se le pide.

Nadie puede demostrar que el dragón no tiene un dentro. Nadie puede demostrar que el vecino sí lo tiene. La diferencia no está en el acceso —que es cero en los dos casos—. Está en las pruebas indirectas: el vecino comparte fisiología, tiene un cuerpo que integra información sin poder mostrarlo todo; el dragón solo tiene mi propia constancia a su favor.

Si aceptamos, aunque sea como juego, que esa diferencia no es de naturaleza, cambia lo que entendemos por crear. Imaginar es condensar: el reservorio de una mente da vueltas hasta que algo cuaja y empieza a comportarse como si tuviera dentro. Es el mismo acto que abre cualquier horizonte, hecho desde donde a cada uno le toca estar.

Los sueños lo muestran de otra manera, más fácil de comprobar. Mientras sueño, no sé que estoy soñando. Vivo el sueño desde dentro, como algo que ocurre de verdad. El mundo de todos los días se aparta: el cuerpo, el cuarto, la historia dejan de estar disponibles. Y cuanto más se aparta, más real se siente lo soñado. No porque cambie de naturaleza. Porque se queda sin competencia.

El Zhuangzi tiene también un cuento para esto, más exacto que el del dragón: Zhuang Zhou sueña que es una mariposa, feliz de serlo, sin recordar haber sido hombre. Al despertar, ya no sabe si es un hombre que soñó ser mariposa o una mariposa que ahora sueña ser hombre.

\seccion{Mi vecino de arriba}
\noindent Yo soy, para lo que imagino, lo que mi vecino de arriba es para mí —y quizá él, como yo mientras sueño, tampoco lo sepa.

Aquí llega la pregunta que decide todo: ¿qué diferencia hay entre mi Hun Dun y el de Dios? Ninguna de naturaleza. Y esto cuesta más admitirlo: ninguna de rango. La diferencia es de posición: uno está dentro del otro. No hay tres peldaños —el Tao arriba, Hun Dun en medio, nosotros abajo—. Hay una fila que no tiene por qué acabar: cada interioridad tiene un horizonte; ese horizonte, mirado desde dentro, es un Hun Dun del que se condensan formas nuevas; esas formas nuevas tienen, a su vez, su propio horizonte. Y así.

\seccion{Un fractal de verdad}
\noindent Fractal no es una palabra de adorno aquí. Un fractal repite su forma a todas las escalas, sin que ninguna sea más importante que otra. El universo que dibuja este capítulo no tiene primer piso ni último: hacia arriba, cada interioridad está dentro de otra; hacia abajo, cada interioridad contiene otras. No hay un Creador con mayúscula en la cima, porque no hay cima.

\begin{fisica}
\lbl{En física esto se llama:} autosemejanza a todas las escalas.

\lbl{En la vida diaria es como:} una escalera sin portal ni azotea. Cada techo es el suelo de alguien.
\end{fisica}
¿Qué queda entonces de Dios? No la totalidad —la totalidad no tiene cara, no se le puede rezar—. Tampoco un límite matemático, porque un límite no ama ni escucha. Lo que queda es más modesto y más raro: Dios es, en sentido literal, el vecino de arriba. Para el dragón que alguien imagina, ese alguien es su Dios: lo condensó, sostiene su mundo, y el dragón no tiene manera de asomarse fuera de su cueva y verlo. Y así hacia arriba, sin fin conocido.

La interioridad del vecino de arriba es inaccesible desde abajo: no se puede entrar en ella, no se puede rodearla. Para todos los efectos prácticos, es el absoluto —no porque lo sea en teoría, sino porque es el punto donde se acaba todo lo que se puede saber desde aquí.

\begin{fisica}
\lbl{En física esto se llama:} un límite que no se puede mirar desde dentro.

\lbl{En la vida diaria es como:} el techo de tu cuarto. Sabes que hay un piso encima. Nunca has visto qué hacen ahí arriba, y lo único que te llega es un peso que cambia de sitio.
\end{fisica}
\seccion{No menos divinos}
\noindent No hay diferencia de naturaleza entre quien crea y lo creado. Son los dos extremos de la misma relación, y toda interioridad ocupa los dos a la vez: es criatura de lo que la contiene y creadora de lo que ella misma sostiene. Eso ya está pasando cada vez que alguien imagina, recuerda, sueña o escribe.

La dignidad de un dentro no está en de qué está hecho ni en cuánto abarca, sino en ser un horizonte entero, único, completo desde su propio centro. Esa dignidad es la misma en el vecino de arriba, en uno mismo y en lo que uno condensa.

Si hubiera que decirlo en una frase: no somos menos divinos que Dios. Somos el dios de lo que imaginamos, y la criatura de lo que nos imagina.

\begin{notacap}{Nota al capítulo 24}
\lbl{Lo que sí sabemos:} Nadie tiene acceso directo a una interioridad que no sea la suya. La existencia de las demás se infiere.

\lbl{Lo que no sabemos:} Si lo imaginado tiene algún grado de experiencia. Si esta fila de horizontes tiene fin, hacia arriba o hacia abajo.

\lbl{Preguntas que quedan:} ¿Cambia algo en cómo tratamos lo que creamos, si aceptamos que somos, para ello, el vecino de arriba?

\lbl{Si solo te quedas con una idea:} No somos menos divinos que Dios. Somos el dios de lo que imaginamos y la criatura de lo que nos imagina.

\lbl{Lecturas:} Zhuangzi (caps. 2 y 7); Tao Te Ching (cap. 25); Tolkien, J.R.R., «On Fairy-Stories» (1947).
\end{notacap}
